% Основные источники:
% https://sochisirius.ru/uploads/2021/11/math_1021_11a_12_gaussovy_celye_chisla_teoriya.pdf
% https://sochisirius.ru/uploads/2021/11/math_1021_11a_13_gaussovy_celye_chisla_zadachi.pdf
% https://sochisirius.ru/uploads/2021/11/math_1021_11b_19-celye_gaussovye_chisla.pdf
% https://sochisirius.ru/uploads/2021/11/math_1021_11v_14_gaussovy_chisla.pdf

\centerline{\bf \Huge Целые гауссовы числа}

\begin{flushright}
        \scriptsize{}
\end{flushright}

\problem{1}{Исследуйте делимость на число $1+i$.}

\begin{enumerate}[label=\textbf{\alph*)}, leftmargin=8mm]
    \item Докажите, что
    \[
        1+i\mid a+bi
        \quad\Longleftrightarrow\quad
        a\equiv b\pmod 2.
    \]

    \item Докажите эквивалентность
    \[
        1+i\mid\alpha
        \quad\Longleftrightarrow\quad
        2\mid N(\alpha).
    \]

    \item Для каждого из чисел
    \[
        8+12i,\qquad 7+i,\qquad 5+3i
    \]
    найдите наибольшее $k$, для которого $(1+i)^k$ делит это число.
\end{enumerate}

\medskip

\problem{2}{Потренируйтесь работать с нормой и делителями.}

\begin{enumerate}[label=\textbf{\alph*)}, leftmargin=8mm]
    \item Найдите все гауссовы числа норм
    \[
        2,\qquad 5,\qquad 13,\qquad 17.
    \]

    \item Разложите числа $2$, $5$, $13$ и $17$ в произведение
    необратимых гауссовых чисел.

    \item Найдите с точностью до ассоциированности все делители числа
    $5$ в кольце $\mathbb Z[i]$.

    \item Докажите, что из $\alpha\mid\beta$, где $\alpha\ne0$, следует
    \[
        N(\alpha)\mid N(\beta).
    \]
    Покажите на примере чисел $1+2i$ и $2+i$, что обратное утверждение
    неверно.
\end{enumerate}

\medskip

\problem{3}{Выполните деление с остатком.}

\begin{enumerate}[label=\textbf{\alph*)}, leftmargin=8mm]
    \item Разделите $17+9i$ на $4+3i$.

    \item Разделите $9+2i$ на $4+7i$.

    \item Найдите два различных допустимых остатка при делении
    $3+i$ на $1+i$.
\end{enumerate}

\begin{taskbox}
    {Задача 4. Рождественская теорема Ферма--Эйлера}
    {red}
    {colback=green!5!white}

Пусть $p$ — обычное простое число вида $4k+1$.

\begin{enumerate}[label=\textbf{\alph*)}, leftmargin=8mm]
    \item Используя теорему Вильсона, докажите, что существует целое
    число $x$, для которого
    \[
        x^2\equiv-1\pmod p.
    \]

    \item В кольце $\mathbb Z[i]$ рассмотрите
    \[
        d=\text{НОД}(p,\;x+i).
    \]
    Докажите, что $d$ не является ни обратимым, ни ассоциированным
    с $p$.

    \item Докажите, что
    \[
        N(d)=p.
    \]

    \item Выведите, что существуют натуральные числа $a,b$, для которых
    \[
        p=a^2+b^2,
        \qquad
        p=(a+bi)(a-bi).
    \]
\end{enumerate}

\end{taskbox}

\medskip

\problem{5}{
Докажите, что с точностью до ассоциированности список всех простых
гауссовых чисел имеет следующий вид:
}

\begin{enumerate}[label=\textbf{\arabic*)}, leftmargin=10mm]
    \item число $1+i$;

    \item обычные простые числа $p\equiv3\pmod4$;

    \item числа $a+bi$, норма которых
    \[
        a^2+b^2
    \]
    является обычным простым числом вида $4k+1$.
\end{enumerate}

\problem{6}{
Докажите, что если простое число $p\equiv1\pmod4$ представлено в виде
\[
    p=a^2+b^2,
\]
то это представление единственно с точностью до перестановки $a,b$
и изменения их знаков.
}

\medskip

\problem{7}{Докажите теорему о сумме двух квадратов.}

\begin{enumerate}[label=\textbf{\alph*)}, leftmargin=8mm]
    \item Получите из равенства
    \[
        (a+bi)(c+di)=(ac-bd)+(ad+bc)i
    \]
    тождество
    \[
        (a^2+b^2)(c^2+d^2)
        =(ac-bd)^2+(ad+bc)^2.
    \]

    \item Докажите, что натуральное число $n$ представимо в виде
    \[
        n=x^2+y^2,\qquad x,y\in\mathbb Z,
    \]
    тогда и только тогда, когда каждый простой делитель
    $p\equiv3\pmod4$ входит в разложение числа $n$ в чётной степени.
\end{enumerate}

\medskip

\problem{8}{
Решите в целых числах уравнение a) $x^2+4=y^3.$ б) $x^2+1=y^3$;
}

% \textit{Указание.} После разбора чётности рассмотрите в
% $\mathbb Z[i]$ равенство
% \[
%     (x+2i)(x-2i)=y^3
% \]
% и исследуйте наибольший общий делитель множителей.

% Источник: листок 11А «Гауссовы целые числа. Задачи», задача 3а.

\medskip

\problem{9}{
Натуральные числа $x,y,z$ удовлетворяют равенству
\[
    xy=z^2+1.
\]
Докажите, что существуют целые числа $a,b,c,d$, для которых
\[
    x=a^2+b^2,\qquad
    y=c^2+d^2,\qquad
    z=ac+bd.
\]
}

% \textit{Указание.} Разложите число $z+i$ на простые множители в
% $\mathbb Z[i]$ и распределите их между двумя множителями с нормами
% $x$ и $y$.

% Источник: листок 11А «Гауссовы целые числа. Задачи», задача 4.

\medskip

\problem{10}{
Найдите все пары простых чисел $p$ и $q$ (не обязательно различных),
для которых каждое из чисел
\[
    2p-1,\qquad 2q-1,\qquad 2pq-1
\]
является квадратом натурального числа.
}

% Источник: листки Сириуса для групп 11А, 11Б и 11В.